% Relations and Logic
\section{Relations and Logical Structure}

\subsection{Binary Relations}

The system implements binary relations $R \subseteq \mathcal{E} \times \mathcal{E}$ where $\mathcal{E}$ is the set of entities.

\subsubsection{Relational Properties}
Observable properties:
\begin{itemize}
    \item \textbf{Non-reflexive}: Entities rarely relate to themselves
    \item \textbf{Asymmetric}: Most relations have clear direction
    \item \textbf{Transitive}: Hierarchical chains suggest transitivity
\end{itemize}

\subsection{Logical Operators}

\subsubsection{Conjunction}
Sequential sign placement suggests conjunction:
\[
\text{ENTITY}_1 \land \text{ENTITY}_2 \land \cdots \land \text{ENTITY}_n
\]

\subsubsection{Absence of Negation}
No evidence for negation operator:
\begin{itemize}
    \item Administrative records state what IS, not what IS NOT
    \item Absence represented by omission, not explicit negation
\end{itemize}

\subsubsection{Conditional Structure}
Possible conditional patterns:
\[
\text{IF } \text{ENTITY}_1 \text{ THEN } \text{ENTITY}_2
\]
represented as sequential validation chains.

\subsection{Propositional Structure}

Inscriptions as propositions:
\begin{itemize}
    \item Each inscription asserts a state or transaction
    \item Truth value: verified (closed/sealed) or unverified (open)
    \item Validation: terminal markers function as truth certification
\end{itemize}

\subsection{Logical Completeness}

The system is NOT logically complete:
\begin{itemize}
    \item Cannot express arbitrary propositions
    \item Limited to specific administrative patterns
    \item Domain-specific rather than general-purpose
\end{itemize}

This is consistent with operational notation rather than logical formalism.

\subsection{Modal Aspects}

Possible modal dimensions:
\begin{itemize}
    \item \textbf{Necessity}: Required vs. optional operations
    \item \textbf{Possibility}: Permitted transaction types
    \item \textbf{Authority}: Authorized vs. unauthorized actions
\end{itemize}

Modal distinctions encoded through sign choice and sequencing.
